\documentclass[11pt,a4paper]{article}

% ------------------------------------------------------------
% Codificacao, idioma e tipografia -- preserva o padrao do Volume I
% ------------------------------------------------------------
\usepackage[utf8]{inputenc}
\usepackage[T1]{fontenc}
\usepackage[brazil]{babel}
\usepackage{lmodern}
\usepackage{microtype}
\usepackage{geometry}
\geometry{top=2.2cm,bottom=2.2cm,left=2.25cm,right=2.25cm}
\usepackage{setspace}
\setstretch{1.07}

% ------------------------------------------------------------
% Matematica e tabelas
% ------------------------------------------------------------
\usepackage{amsmath,amssymb,mathtools}
\usepackage{bm}
\usepackage{dsfont}
\usepackage{array,booktabs,longtable,tabularx,multirow}
\usepackage{enumitem}
\usepackage{siunitx}
\sisetup{output-decimal-marker={,},group-separator={.}}

% ------------------------------------------------------------
% Graficos e caixas -- mesma paleta do Volume I
% ------------------------------------------------------------
\usepackage{xcolor}
\definecolor{Navy}{HTML}{18324A}
\definecolor{Blue}{HTML}{2D628F}
\definecolor{Teal}{HTML}{2D7A78}
\definecolor{Gold}{HTML}{B8862C}
\definecolor{Red}{HTML}{9B3A36}
\definecolor{LightBlue}{HTML}{EDF4F9}
\definecolor{LightTeal}{HTML}{EEF7F5}
\definecolor{LightGold}{HTML}{FBF6EA}
\definecolor{LightRed}{HTML}{FAEFEE}
\definecolor{LightGray}{HTML}{F4F5F6}
\definecolor{MidGray}{HTML}{7A838A}

\usepackage{tikz}
\usetikzlibrary{arrows.meta,positioning,calc,fit,backgrounds,shapes.geometric,matrix}
\usepackage[most]{tcolorbox}

\newtcolorbox{intuicao}[1][]{
  enhanced,breakable,colback=LightBlue,colframe=Blue,boxrule=0.6pt,
  arc=1.5mm,left=2mm,right=2mm,top=1.5mm,bottom=1.5mm,
  title=Intuição,#1}
\newtcolorbox{formalizacao}[1][]{
  enhanced,breakable,colback=LightTeal,colframe=Teal,boxrule=0.6pt,
  arc=1.5mm,left=2mm,right=2mm,top=1.5mm,bottom=1.5mm,
  title=Formalização,#1}
\newtcolorbox{cuidado}[1][]{
  enhanced,breakable,colback=LightRed,colframe=Red,boxrule=0.6pt,
  arc=1.5mm,left=2mm,right=2mm,top=1.5mm,bottom=1.5mm,
  title=Precisão conceitual,#1}
\newtcolorbox{fonte}[1][]{
  enhanced,breakable,colback=LightGold,colframe=Gold,boxrule=0.6pt,
  arc=1.5mm,left=2mm,right=2mm,top=1.5mm,bottom=1.5mm,
  title=Base normativa,#1}

% ------------------------------------------------------------
% Cabecalho, hiperlinks, secoes
% ------------------------------------------------------------
\usepackage{fancyhdr}
\pagestyle{fancy}
\fancyhf{}
\fancyhead[L]{\small\color{MidGray}Parametrização matemática -- Volume II}
\fancyhead[R]{\small\color{MidGray}\thepage}
\renewcommand{\headrulewidth}{0.3pt}
\setlength{\headheight}{14pt}

\usepackage{titlesec}
\titleformat{\section}{\Large\bfseries\color{Navy}}{\thesection}{0.75em}{}
\titleformat{\subsection}{\large\bfseries\color{Blue}}{\thesubsection}{0.65em}{}
\titleformat{\subsubsection}{\normalsize\bfseries\color{Teal}}{\thesubsubsection}{0.55em}{}
\titlespacing*{\section}{0pt}{1.8em}{0.7em}
\titlespacing*{\subsection}{0pt}{1.35em}{0.5em}

\usepackage{xurl}
\usepackage[hidelinks]{hyperref}
\hypersetup{
  pdftitle={Parametrizacao matematica de objetos contabeis e tributarios - Volume II},
  pdfsubject={Regimes tributarios, dados minimos, contrafactuais e rastreabilidade normativa},
  pdfauthor={Caderno tecnico},
  colorlinks=true,
  linkcolor=Navy,
  urlcolor=Blue,
  citecolor=Teal
}
\usepackage[nameinlink,noabbrev]{cleveref}

% ------------------------------------------------------------
% Comandos coerentes com o Volume I
% ------------------------------------------------------------
\newcommand{\R}{\mathbb{R}}
\newcommand{\N}{\mathbb{N}}
\newcommand{\ind}{\mathds{1}}
\newcommand{\BP}{\mathrm{BP}}
\newcommand{\DRE}{\mathrm{DRE}}
\newcommand{\DFC}{\mathrm{DFC}}
\newcommand{\DVA}{\mathrm{DVA}}
\newcommand{\CBS}{\mathrm{CBS}}
\newcommand{\IBS}{\mathrm{IBS}}
\newcommand{\IRPJ}{\mathrm{IRPJ}}
\newcommand{\CSLL}{\mathrm{CSLL}}
\newcommand{\cT}{\mathcal{T}}
\newcommand{\cB}{\mathcal{B}}
\newcommand{\cG}{\mathcal{G}}
\newcommand{\cA}{\mathcal{A}}
\newcommand{\cC}{\mathcal{C}}
\newcommand{\cD}{\mathcal{D}}
\newcommand{\cH}{\mathcal{H}}
\newcommand{\cS}{\mathcal{S}}
\newcommand{\cQ}{\mathcal{Q}}
\newcommand{\cN}{\mathcal{N}}
\newcommand{\periodo}[1]{I_{#1}:=(#1,#1+1]}
\newcommand{\Prov}{\operatorname{Prov}}
\newcommand{\Dep}{\operatorname{Dep}}
\newcommand{\Adm}{\operatorname{Adm}}
\newcommand{\TaxExp}{E^{\mathrm{trib}}}
\newcommand{\TaxCash}{P^{\mathrm{tax,cash}}}

\setlist[itemize]{leftmargin=1.4em,itemsep=0.25em,topsep=0.35em}
\setlist[enumerate]{leftmargin=1.7em,itemsep=0.3em,topsep=0.4em}

\begin{document}
\pagenumbering{gobble}

% ============================================================
% CAPA
% ============================================================
\begin{titlepage}
\centering
\vspace*{1.2cm}
{\Large\color{Blue}\textsc{Caderno de formalização -- Volume II}}\\[0.50cm]
{\huge\bfseries\color{Navy}Regime tributário, dados mínimos\\[0.15cm]
e análise contrafactual}\\[0.65cm]
{\Large\color{Teal}Extensão da parametrização matemática de objetos\\
contábeis e tributários}\\[1.15cm]

\begin{tcolorbox}[width=0.90\textwidth,colback=LightGray,colframe=MidGray,boxrule=0.5pt,arc=2mm]
\centering
\large
\[
\boxed{\zeta_t=\bigl(x_t^{\min},u_t^{\min},\rho_t,\eta_t\bigr)}
\]
\vspace{-1mm}
\small
Dados empresariais observados $\to$ configuração tributária aplicável $\to$\\
operadores efetivos $\to$ cenário comparável e auditável.
\end{tcolorbox}

\vfill
{\large Versão de referência: 1º de setembro de 2026}\\[0.15cm]
{\small Documento complementar ao Volume I; não constitui parecer contábil ou jurídico.}
\end{titlepage}

\pagenumbering{roman}
\tableofcontents
\clearpage
\pagenumbering{arabic}

% ============================================================
\section{Escopo diferencial em relação ao Volume I}
% ============================================================

O Volume I já formalizou o núcleo do sistema: estado $x_t$, eventos $u_t$, regras contábeis e tributárias, demonstrações, operadores $\cB_j$, $\cA_j$, $\cC_j$ e $\cD_j$, apuração, hierarquia normativa e um primeiro experimento contrafactual. Este Volume II trata apenas da camada que apareceu posteriormente na discussão e que não estava parametrizada, ou estava apenas mencionada.

\begin{formalizacao}
O diferencial pode ser resumido em cinco objetos novos ou refinados:
\[
\boxed{
\rho_t,\qquad \eta_t,\qquad \zeta_t,\qquad \mathfrak E_{j,t},\qquad \Prov.
}
\]
Eles representam, respectivamente, a configuração tributária da entidade, seu perfil factual, o vetor mínimo de entrada, o seletor de regras efetivamente aplicáveis e a proveniência normativa de cada parâmetro.
\end{formalizacao}

\subsection{Correções de notação necessárias}

No Volume I, $z_t$ já foi utilizado como símbolo genérico para ajustes e informações fiscais adicionais na construção de uma base tributável. A discussão posterior reutilizou $z_t$ para o vetor completo de entrada do simulador. Para evitar colisão semântica, este Volume II preserva $z_t$ no sentido original e introduz
\begin{equation}
\boxed{\zeta_t:=\bigl(x_t^{\min},u_t^{\min},\rho_t,\eta_t\bigr).}
\label{eq:zeta}
\end{equation}

Também preservaremos $s_{k,t}$, já usado no Volume I para data/instante da operação. O eventual tratamento especial de uma operação será denotado por $h_{k,t}$, e não por $s_{k,t}$.

\begin{cuidado}
A comparação realizada para este volume mostrou que o Volume I já mencionava razão/diário, documentos fiscais e uma camada de lançamentos $\Lambda_t$, além de já conter um operador contrafactual. Portanto, esses temas não serão reapresentados do zero. O avanço aqui consiste em parametrizar \emph{quais dados mínimos são necessários}, \emph{como o regime da entidade entra no operador tributário} e \emph{como cada regra fica rastreável até sua fonte oficial}.
\end{cuidado}

% ============================================================
\section{O regime tributário como objeto independente $\rho_t$}
% ============================================================

\subsection{Lei, regime da entidade e fatos da entidade não são o mesmo objeto}

No Volume I, $\Theta_t^{\mathrm{tax}}=(\theta_{j,t})_{j\in\mathsf J_t}$ representa a legislação e regulamentação tributária vigente. Isso, por si só, não especifica qual ramo dessa legislação se aplica a uma entidade concreta.

Introduzimos então
\begin{equation}
\boxed{\rho_t\in\mathsf R_t,}
\label{eq:rho}
\end{equation}
onde $\rho_t$ codifica as escolhas, enquadramentos e regimes de apuração relevantes para a entidade no período $t$.

\begin{intuicao}
$\Theta_t^{\mathrm{tax}}$ responde: ``quais são as regras disponíveis no sistema jurídico?''. Já $\rho_t$ responde: ``em qual configuração tributária esta entidade se encontra?''. Um terceiro objeto, $\eta_t$, descreverá fatos da própria entidade que condicionam elegibilidade e aplicação das regras.
\end{intuicao}

\subsection{Decomposição por eixos}

Uma parametrização mais fiel do que tratar $\rho_t$ como um rótulo único é
\begin{equation}
\boxed{
\rho_t=
\bigl(
\rho_t^{\mathrm{ent}},
\rho_t^{\mathrm{IR}},
\rho_t^{\mathrm{cons}},
\rho_t^{\mathrm{esp}}
\bigr).
}
\label{eq:rho-decomp}
\end{equation}

Os componentes têm interpretação distinta:
\begin{itemize}
    \item $\rho_t^{\mathrm{ent}}$: enquadramento geral da entidade, por exemplo opção pelo Simples Nacional ou situação fora do Simples;
    \item $\rho_t^{\mathrm{IR}}$: regime/base de apuração aplicável a IRPJ/CSLL quando pertinente, como lucro real, presumido ou arbitrado;
    \item $\rho_t^{\mathrm{cons}}$: regime de apuração de tributos sobre consumo, incluindo a posição da entidade perante o regime regular de IBS/CBS;
    \item $\rho_t^{\mathrm{esp}}$: regimes ou tratamentos específicos adicionais relevantes ao caso.
\end{itemize}

\begin{fonte}
A Receita Federal informa, na documentação da ECF de 2026, que pessoas jurídicas podem estar tributadas pelo lucro real, arbitrado ou presumido, com exceção dos optantes do Simples Nacional em relação à obrigação de ECF. A LC 214/2025, art. 41, distingue o regime regular de IBS/CBS das regras do Simples Nacional/MEI e prevê que optantes do Simples possam optar por apurar e recolher IBS/CBS pelo regime regular. Ver \cite{RFB_ECF_2026,LC214}.
\end{fonte}

\begin{cuidado}
O Simples Nacional não deve ser modelado simplesmente como ``mais uma forma de lucro'' ao lado de real, presumido e arbitrado. Ele é um regime especial e unificado com lógica própria. A decomposição em eixos evita misturar natureza jurídica, base de IRPJ/CSLL e regime de IBS/CBS.
\end{cuidado}

% ============================================================
\section{Perfil factual da entidade: $\eta_t$}
% ============================================================

Certas regras dependem não apenas da opção tributária, mas de fatos da entidade e da operação. Introduzimos
\begin{equation}
\boxed{
\eta_t=
(\eta_{1,t},\ldots,\eta_{r,t})
}
\label{eq:eta}
\end{equation}
como vetor de atributos de contexto relevantes ao modelo.

Uma decomposição ilustrativa é
\begin{equation}
\eta_t=
(\text{atividade},\text{localização},\text{porte},\text{cadastros},
\text{benefícios},\text{condições especiais},\ldots).
\end{equation}

\begin{cuidado}
Esses campos não são universalmente obrigatórios em qualquer simulador. A regra metodológica é incluir em $\eta_t$ apenas atributos dos quais algum operador tributário ou teste de elegibilidade realmente dependa. ``Setor'', ``porte'' ou ``localização'' são exemplos, não uma lista normativa fechada.
\end{cuidado}

\subsection{Regime admissível}

Nem todo $\rho\in\mathsf R_t$ é legalmente admissível para qualquer entidade. Definimos uma função de admissibilidade
\begin{equation}
\boxed{
\chi_t(\rho;\eta_t,\Theta_t^{\mathrm{tax}})\in\{0,1\}
}
\label{eq:chi}
\end{equation}
e o conjunto de cenários admissíveis
\begin{equation}
\boxed{
\mathsf R_t^{\mathrm{adm}}(\eta_t;\Theta_t^{\mathrm{tax}})
:=\{\rho\in\mathsf R_t:\chi_t(\rho;\eta_t,\Theta_t^{\mathrm{tax}})=1\}.
}
\label{eq:r-adm}
\end{equation}

\begin{intuicao}
Antes de perguntar ``quanto a empresa pagaria no regime $\rho$?'', o simulador deve perguntar ``a empresa pode estar nesse regime?''. Isso impede comparar cenários matematicamente definidos, mas juridicamente impossíveis.
\end{intuicao}

% ============================================================
\section{Da norma geral à regra efetiva: o seletor $\mathfrak E_{j,t}$}
% ============================================================

O Volume I aplicava diretamente $\theta_{j,t}$ às operações. Com $\rho_t$ e $\eta_t$ explícitos, é útil introduzir um estágio intermediário:
\begin{equation}
\boxed{
\theta_{j,t}^{\mathrm{eff}}
=\mathfrak E_{j,t}
\bigl(\theta_{j,t},\rho_t,\eta_t\bigr).
}
\label{eq:effective-rule}
\end{equation}

$\mathfrak E_{j,t}$ não altera a lei. Ele seleciona, entre as regras vigentes, quais disposições são efetivamente aplicáveis ao contribuinte e à situação considerada.

A cadeia tributária passa a ser
\begin{equation}
(\cT_{k,t},\rho_t,\eta_t)
\xrightarrow{\theta_{j,t}}
\theta_{j,t}^{\mathrm{eff}}
\xrightarrow{\cB_j,\cA_j,\cC_j,\cD_j}
(B_{j,k,t},\tau_{j,k,t},C_{j,k,t},D_{j,k,t}).
\label{eq:pipeline-effective}
\end{equation}

\begin{formalizacao}
No nível do sistema,
\[
\Theta_t^{\mathrm{eff}}
:=\bigl(\mathfrak E_{j,t}(\theta_{j,t},\rho_t,\eta_t)\bigr)_{j\in\mathsf J_t}.
\]
O operador tributário do simulador deve depender de $\Theta_t^{\mathrm{eff}}$, não apenas do texto normativo bruto.
\end{formalizacao}

\begin{fonte}
A necessidade dessa seleção é observável, por exemplo, no art. 41 da LC 214/2025: a mesma legislação contém ramos diferentes para contribuinte sujeito ao regime regular, optante do Simples/MEI e optante do Simples que escolhe apurar IBS/CBS pelo regime regular. Ver \cite{LC214,RFB_SIMPLES_2026}.
\end{fonte}

% ============================================================
\section{Vetor de entrada mínimo como problema de suficiência}
% ============================================================

\subsection{Definição do vetor mínimo}

O vetor complementar proposto na conversa passa a ser, com notação corrigida,
\begin{equation}
\boxed{
\zeta_t=
\left(
 x_t^{\min},
 u_t^{\min},
 \rho_t,
 \eta_t
\right).
}
\end{equation}

A palavra ``mínimo'' não significa ``menor número possível de colunas'' de forma absoluta. Significa: informação suficiente para reproduzir os resultados tributários de interesse no escopo definido.

\subsection{Projeções mínimas}

Se $x_t$ é o estado completo e $u_t$ a coleção granular completa, escrevemos
\begin{equation}
\boxed{
 x_t^{\min}=\Pi_x^{\mathcal H}(x_t),
 \qquad
 u_t^{\min}=\Pi_u^{\mathcal H}(u_t),
}
\label{eq:min-projections}
\end{equation}
em que as projeções retêm apenas os componentes consumidos pelo operador $\cH$.

Podemos caracterizar os campos necessários por
\begin{equation}
\Dep(\cH):=\{\text{atributos lidos por }\cH\}.
\end{equation}

\begin{formalizacao}
Uma projeção $\Pi$ da informação operacional é \textbf{suficiente para o simulador} se existe $\widetilde{\cH}$ tal que
\[
\cH(u_t)=\widetilde{\cH}(\Pi(u_t)).
\]
Equivalentemente,
\[
\Pi(u_t)=\Pi(u_t')
\quad\Longrightarrow\quad
\cH(u_t)=\cH(u_t').
\]
Essa é uma noção de suficiência funcional do dado para o cálculo, e não uma definição estatística de suficiência de Fisher.
\end{formalizacao}

\subsection{Componentes de estado que podem ser necessários}

Em vez de fixar uma lista universal, tratamos
\begin{equation}
 x_t^{\min}
 =\bigl(
 \text{saldos de estado exigidos pelos operadores vigentes}
 \bigr).
\end{equation}

Exemplos possíveis incluem saldos de créditos tributários, obrigações tributárias, prejuízos/bases fiscais carregados, estoques ou outros estados auxiliares quando uma regra possui memória intertemporal.

\begin{fonte}
A documentação oficial da ECF e o CPC 32 evidenciam que a apuração de tributos sobre o lucro pode exigir bases fiscais, prejuízos fiscais, créditos e diferenças temporárias que conectam períodos. Ver \cite{RFB_ECF_2026,CPC32}.
\end{fonte}

\subsection{Operação mínima compatível com a notação do Volume I}

O Volume I definiu
\[
\cT_{k,t}=(d_{k,t},q_{k,t},v_{k,t},s_{k,t},a_{k,t},p_{k,t},\ell_{k,t},\ldots).
\]
Para o problema tributário, refinamos $\ell_{k,t}$ e adicionamos localização e tratamento especial:
\begin{equation}
\boxed{
\cT_{k,t}^{\min}=
(d_{k,t},q_{k,t},v_{k,t},s_{k,t},\ell_{k,t}^{\min},g_{k,t},h_{k,t}),
}
\label{eq:tmin}
\end{equation}
onde:
\begin{itemize}
  \item $d_{k,t}$: direção entrada/saída;
  \item $q_{k,t}$: natureza econômica relevante;
  \item $v_{k,t}$: valor monetário bruto;
  \item $s_{k,t}$: data/instante da operação, preservando a notação do Volume I;
  \item $\ell_{k,t}^{\min}$: códigos/classificações contábeis e fiscais efetivamente necessários;
  \item $g_{k,t}$: atributos de origem, destino ou local da operação quando relevantes;
  \item $h_{k,t}$: marcador de tratamento específico/diferenciado quando exigido.
\end{itemize}

\begin{cuidado}
O conteúdo exato de $\ell^{\min}$, $g$ e $h$ depende do tributo e da versão normativa. A definição correta deve ser obtida por dependência dos operadores, e não por uma lista fixa inventada pelo modelador.
\end{cuidado}

% ============================================================
\section{Por que plano de contas ou balancete, isoladamente, podem ser insuficientes}
% ============================================================

Seja
\begin{equation}
\mathcal A_{\mathrm{acct}}:u_t\longmapsto b_t
\end{equation}
um operador que agrega a atividade do período em um balancete ou conjunto de saldos contábeis $b_t$.

Se existirem duas coleções de operações $u_t$ e $u_t'$ tais que
\begin{equation}
\mathcal A_{\mathrm{acct}}(u_t)=\mathcal A_{\mathrm{acct}}(u_t')
\quad\text{mas}\quad
\cH^{\mathrm{tax}}(u_t)\neq\cH^{\mathrm{tax}}(u_t'),
\label{eq:insuf-aggregate}
\end{equation}
então o resultado tributário não pode ser reconstruído unicamente a partir de $b_t$.

\begin{intuicao}
Duas empresas podem exibir a mesma receita total em um balancete e, ainda assim, ter operações com classificações fiscais, locais, destinatários ou tratamentos diferentes. Se alguma dessas diferenças altera $\cB_j$, $\cA_j$ ou $\cC_j$, a agregação contábil perdeu informação necessária ao recálculo tributário.
\end{intuicao}

Definimos o conjunto de resultados compatíveis com um agregado $b_t$ por
\begin{equation}
\boxed{
\mathcal Y(b_t)
:=\{\cH^{\mathrm{tax}}(u):\mathcal A_{\mathrm{acct}}(u)=b_t\}.
}
\end{equation}
Se $\mathcal Y(b_t)$ não é unitário, existe \emph{não identificação} do resultado tributário a partir daquele agregado.

\begin{fonte}
A orientação oficial da RTC para 2026 exige, para os documentos alcançados, destaque de CBS e IBS individualizado por operação segundo os leiautes técnicos. Esse nível operacional é evidência direta de que a camada tributária trabalha com atributos mais granulares do que um total contábil agregado. Ver \cite{RFB_ORIENT_2026,RFB_RTC_LEG}.
\end{fonte}

% ============================================================
\section{Camada prática de dados: de registros reais a $\zeta_t$}
% ============================================================

\subsection{Pacote de fontes empresariais}

Para aproximar a abstração dos arquivos efetivamente disponíveis, definimos
\begin{equation}
\boxed{
\cQ_t=
(\mathsf{ECD}_t,\mathsf{ECF}_t,\mathsf{DocFisc}_t,\mathsf{Cad}_t,\mathsf{Aux}_t),
}
\label{eq:qdata}
\end{equation}
onde os componentes representam, quando aplicáveis:
\begin{itemize}
  \item $\mathsf{ECD}_t$: escrituração contábil digital e estruturas associadas, como plano de contas, lançamentos e saldos;
  \item $\mathsf{ECF}_t$: escrituração contábil-fiscal e controles relativos à tributação da renda;
  \item $\mathsf{DocFisc}_t$: documentos fiscais eletrônicos e seus atributos por operação;
  \item $\mathsf{Cad}_t$: dados cadastrais e de enquadramento;
  \item $\mathsf{Aux}_t$: módulos auxiliares necessários, como estoques, contas a receber/pagar ou controles fiscais específicos.
\end{itemize}

A reconstrução do vetor de entrada é um operador de integração de dados:
\begin{equation}
\boxed{
\widehat\zeta_t=\cN_t(\cQ_t),
}
\label{eq:normalize}
\end{equation}
em que $\cN_t$ executa extração, normalização, conciliação e mapeamento para o esquema formal.

\begin{fonte}
O Manual da ECD 2026 descreve a escrituração digital e registros de plano de contas, saldos e lançamentos; a Receita Federal mantém documentação própria da ECF para as bases fiscais; e as orientações da RTC de 2026 especificam documentos fiscais eletrônicos com IBS/CBS individualizados por operação. Ver \cite{ECD2026,RFB_ECF_2026,RFB_ORIENT_2026}.
\end{fonte}

\subsection{Relação com a camada $\Lambda_t$ do Volume I}

A camada de lançamentos/razão $\Lambda_t$ do Volume I não é substituída. Ela pode ser recuperada a partir de $\mathsf{ECD}_t$ e sistemas contábeis, enquanto $\mathsf{DocFisc}_t$ fornece atributos transacionais que podem não sobreviver integralmente à agregação contábil.

\begin{formalizacao}
Uma arquitetura de integração compatível com os dois volumes é
\[
\cQ_t
\xrightarrow{\cN_t}
\zeta_t
\xrightarrow{\mathfrak E_t}
\Theta_t^{\mathrm{eff}}
\xrightarrow{\cH}
Y_t.
\]
\end{formalizacao}

\begin{center}
\begin{tikzpicture}[
  node distance=8mm and 8mm,
  box/.style={draw=Navy,rounded corners=2mm,fill=LightBlue,align=center,minimum height=10mm,minimum width=27mm,font=\small},
  op/.style={draw=Teal,rounded corners=2mm,fill=LightTeal,align=center,minimum height=10mm,minimum width=26mm,font=\small},
  arr/.style={-{Latex[length=2.2mm]},thick,Navy}
]
\node[box] (q) {$\cQ_t$\\fontes reais};
\node[op,right=of q] (n) {$\cN_t$\\normalização};
\node[box,right=of n] (z) {$\zeta_t$\\entrada mínima};
\node[op,right=of z] (e) {$\mathfrak E_t$\\regras efetivas};
\node[box,right=of e] (y) {$Y_t$\\saídas};
\draw[arr] (q)--(n);
\draw[arr] (n)--(z);
\draw[arr] (z)--(e);
\draw[arr] (e)--(y);
\end{tikzpicture}
\end{center}

% ============================================================
\section{Experimento contrafactual com lei e regime}
% ============================================================

O Volume I comparava $\Theta^{\mathrm{old}}$ e $\Theta^{\mathrm{new}}$ mantendo $(x_t,u_t)$ fixos. A nova camada permite separar três efeitos: mudança legislativa, mudança de regime e interação entre ambos.

Defina a base econômica observada sem o regime:
\begin{equation}
\boxed{
\bar\zeta_t:=(x_t^{\min},u_t^{\min},\eta_t).
}
\end{equation}
Para um cenário $s$,
\begin{equation}
\zeta_t^{(s)}=(\bar\zeta_t,\rho_t^{(s)}),
\qquad
Y_t^{(s)}=\cH(\zeta_t^{(s)};\Theta_t^{(s)}).
\end{equation}

\subsection{Decomposição dos efeitos}

Com cenário inicial $(\rho^0,\Theta^0)$ e cenário final $(\rho^1,\Theta^1)$, definimos
\begin{align}
\Delta_{\Theta}Y_t
&:=\cH(\bar\zeta_t,\rho^0;\Theta^1)
   -\cH(\bar\zeta_t,\rho^0;\Theta^0),\\
\Delta_{\rho}Y_t
&:=\cH(\bar\zeta_t,\rho^1;\Theta^0)
   -\cH(\bar\zeta_t,\rho^0;\Theta^0),\\
\Delta_{\mathrm{joint}}Y_t
&:=\cH(\bar\zeta_t,\rho^1;\Theta^1)
   -\cH(\bar\zeta_t,\rho^0;\Theta^0).
\end{align}

A interação residual é
\begin{equation}
\boxed{
\Delta_{\mathrm{int}}Y_t
:=\Delta_{\mathrm{joint}}Y_t
 -\Delta_{\Theta}Y_t
 -\Delta_{\rho}Y_t.
}
\label{eq:interaction}
\end{equation}

\begin{intuicao}
Essa decomposição evita atribuir à ``Reforma'' um efeito que, na verdade, veio de uma troca de regime da empresa; e evita atribuir à troca de regime um efeito causado pela nova lei. Quando ambos mudam, $\Delta_{\mathrm{int}}$ mede a parte que não é explicada pela soma dos efeitos isolados.
\end{intuicao}

\begin{cuidado}
Cada cenário envolvendo $\rho^s$ deve satisfazer $\rho^s\in\mathsf R_t^{\mathrm{adm}}(\eta_t;\Theta^s)$. O contrafactual é condicional à manutenção de $\bar\zeta_t$; ele isola o efeito normativo/paramétrico, mas não modela reações de preço, quantidade, investimento ou reorganização operacional.
\end{cuidado}

% ============================================================
\section{Qual impacto medir? Apuração, caixa e despesa tributária}
% ============================================================

A conversa posterior propôs comparar ``despesas com impostos''. Esse objetivo exige uma distinção adicional: montante apurado, valor a recolher, pagamento de caixa e despesa contábil não são o mesmo objeto.

Definimos, por tributo $j$:
\begin{equation}
\boxed{
Y_{j,t}^{\mathrm{tax}}=
\left(
S_{j,t}^{\mathrm{apur}},
T_{j,t}^{\mathrm{recolher}},
P_{j,t}^{\mathrm{cash}},
E_{j,t}^{\mathrm{DRE}},
C_{j,t}^{\mathrm{saldo}}
\right).
}
\label{eq:tax-output}
\end{equation}

Aqui:
\begin{itemize}
  \item $S_{j,t}^{\mathrm{apur}}$: saldo de apuração, já definido no Volume I;
  \item $T_{j,t}^{\mathrm{recolher}}$: obrigação positiva a recolher;
  \item $P_{j,t}^{\mathrm{cash}}$: pagamento financeiro realizado no período;
  \item $E_{j,t}^{\mathrm{DRE}}$: efeito reconhecido como despesa/resultado segundo a norma contábil aplicável;
  \item $C_{j,t}^{\mathrm{saldo}}$: saldo credor/ativo recuperável quando aplicável.
\end{itemize}

\subsection{Tributos sobre o lucro}

Para tributos sobre o lucro, o CPC 32 distingue tributo corrente e diferido. Em notação resumida,
\begin{equation}
\boxed{
E_t^{\mathrm{IR,trib}}
=E_t^{\mathrm{corrente}}+E_t^{\mathrm{diferido}}
}
\end{equation}
no escopo da despesa/receita tributária tratada pelo pronunciamento.

\subsection{Tributos recuperáveis sobre aquisições}

Para estoques, o CPC 16 determina que tributos recuperáveis junto ao fisco não compõem o custo de aquisição. Portanto, a existência de um crédito tributário recuperável impede a identificação ingênua
\[
\text{tributo destacado ou apurado}\equiv\text{despesa contábil}.
\]

\begin{fonte}
O CPC 32 define despesa tributária como composta por despesa tributária corrente e diferida no contexto de tributos sobre o lucro. O CPC 16 exclui do custo de aquisição de estoques os tributos recuperáveis junto ao fisco. Ver \cite{CPC32,CPC16}.
\end{fonte}

\begin{formalizacao}
Para o projeto, a função-objetivo deve ser declarada explicitamente. Exemplos:
\[
\Delta T_t^{\mathrm{recolher}},\qquad
\Delta P_t^{\mathrm{cash}},\qquad
\Delta E_t^{\mathrm{DRE}},\qquad
\Delta C_t^{\mathrm{saldo}},\qquad
\Delta L_t^{\mathrm{ap\acute{o}s\ tributos}}.
\]
A expressão ``impacto tributário'' é ambígua enquanto esse vetor de resposta não estiver fixado.
\end{formalizacao}

% ============================================================
\section{Camada de proveniência e versionamento normativo}
% ============================================================

O Volume I já apresentou a hierarquia constitucional, legal e regulamentar. O avanço necessário para um simulador auditável é associar cada parâmetro implementado a uma fonte e a uma versão.

\subsection{Taxonomia de fontes}

Definimos
\begin{equation}
\boxed{
\cS=
\{S^{\mathrm{norm}},S^{\mathrm{reg}},S^{\mathrm{tec}},S^{\mathrm{oper}}\},
}
\label{eq:sources}
\end{equation}
com classificação operacional:
\begin{itemize}
  \item $S^{\mathrm{norm}}$: Constituição, emendas, leis complementares e demais normas legais de referência;
  \item $S^{\mathrm{reg}}$: decretos, resoluções, instruções e atos regulamentares;
  \item $S^{\mathrm{tec}}$: estudos e materiais técnicos oficiais, úteis para interpretação e parametrização, sem substituir a norma;
  \item $S^{\mathrm{oper}}$: manuais, leiautes, notas técnicas e documentação de sistemas/documentos eletrônicos.
\end{itemize}

\begin{cuidado}
$\cS$ é uma classificação de fontes para engenharia do modelo, não uma nova hierarquia jurídica. Documentos técnicos e operacionais não prevalecem sobre normas de nível superior.
\end{cuidado}

\subsection{Proveniência de cada parâmetro}

Para cada parâmetro, regra ou ramificação implementada $p$, armazenamos
\begin{equation}
\boxed{
\Prov(p)=
(\text{fonte},\text{dispositivo},\text{versão},\text{vigência},\text{data de consulta}).
}
\label{eq:provenance}
\end{equation}

A regra metodológica é
\begin{equation}
\boxed{
\forall p\in\Theta_t^{\mathrm{eff}},\qquad \Prov(p)\neq\varnothing.
}
\label{eq:trace-rule}
\end{equation}

Também convém versionar a função implementada:
\begin{equation}
\cH_t=\cH^{[\nu(t)]},
\end{equation}
em que $\nu(t)$ identifica a versão normativa e computacional usada no período.

\begin{fonte}
A Receita Federal mantém página consolidada da legislação da RTC com EC 132/2023, LC 214/2025, LC 227/2026, Decreto 12.955/2026, atos conjuntos e atos técnicos. A mesma instituição, em parceria com o CFC, disponibiliza curso oficial cujo Módulo 1 cobre fato gerador, base de cálculo, alíquotas, local da operação, não cumulatividade e apuração. Ver \cite{RFB_RTC_LEG,RFB_CURSO}.
\end{fonte}

% ============================================================
\section{Exemplos}
% ============================================================

\subsection{Exemplo ilustrativo I -- por que um agregado pode não identificar o tributo}

Suponha que um balancete informe apenas receita total de R\$ 200. Existem duas decomposições compatíveis:
\[
u_t^{(A)}=(\cT_{1,t}^{(A)},\cT_{2,t}^{(A)}),
\qquad
u_t^{(B)}=(\cT_{1,t}^{(B)},\cT_{2,t}^{(B)}),
\]
com
\[
\sum_k v_{k,t}^{(A)}=\sum_k v_{k,t}^{(B)}=200.
\]
Se as classificações fiscais ou locais das operações diferirem e isso alterar $\cA_j$ ou $\cB_j$, então
\[
\mathcal A_{\mathrm{acct}}(u_t^{(A)})
=\mathcal A_{\mathrm{acct}}(u_t^{(B)})
\quad\text{e}\quad
\cH^{\mathrm{tax}}(u_t^{(A)})\neq\cH^{\mathrm{tax}}(u_t^{(B)}).
\]
Logo, o balancete é insuficiente para esse cálculo. Não é necessário supor números de alíquotas: o resultado é estrutural.

\subsection{Exemplo ilustrativo II -- separação dos efeitos de lei e regime}

Considere um resultado escalar $Y$ e valores hipotéticos:
\[
\cH(\rho^0;\Theta^0)=100,
\quad
\cH(\rho^0;\Theta^1)=108,
\quad
\cH(\rho^1;\Theta^0)=96,
\quad
\cH(\rho^1;\Theta^1)=101.
\]
Então
\begin{align*}
\Delta_{\Theta}Y&=108-100=8,\\
\Delta_{\rho}Y&=96-100=-4,\\
\Delta_{\mathrm{joint}}Y&=101-100=1,\\
\Delta_{\mathrm{int}}Y&=1-8-(-4)=-3.
\end{align*}
A troca conjunta produz aumento de 1, mas isso não significa que ``a nova lei aumentou 1''. A decomposição revela efeitos opostos e uma interação de $-3$.

\subsection{Exemplo legal I -- Simples Nacional e regime regular de IBS/CBS}

A LC 214/2025, art. 41, permite representar uma bifurcação real do seletor de regras. Para uma entidade optante do Simples Nacional:
\[
\rho_t^{\mathrm{cons}}
\in
\{\text{regras do Simples},\text{opção pelo regime regular de IBS/CBS}\},
\]
observadas as condições legais.

No modelo:
\[
\theta_{\IBS,t}^{\mathrm{eff}}
=\mathfrak E_{\IBS,t}(\theta_{\IBS,t},\rho_t,\eta_t),
\qquad
\theta_{\CBS,t}^{\mathrm{eff}}
=\mathfrak E_{\CBS,t}(\theta_{\CBS,t},\rho_t,\eta_t).
\]
O exemplo é de \emph{seleção de ramo normativo}; não atribuímos aqui alíquota ou vantagem econômica sem os dados concretos da empresa.

\begin{fonte}
O art. 41 da LC 214/2025 disciplina o regime regular e a opção de contribuinte do Simples pelo regime regular. Em agosto de 2026, o CGSN atualizou regras do Simples para adaptação à RTC, incluindo IBS e CBS. Ver \cite{LC214,RFB_SIMPLES_2026}.
\end{fonte}

\subsection{Exemplo legal-operacional II -- granularidade documental em 2026}

A orientação oficial para 2026 determina que os documentos fiscais eletrônicos alcançados sejam emitidos com destaque de CBS e IBS individualizados por operação conforme notas técnicas. Logo, uma implementação pode usar
\[
\mathsf{DocFisc}_t\xrightarrow{\cN_t}u_t^{\min}
\]
como uma das rotas de reconstrução da camada granular.

\begin{fonte}
A Receita Federal lista, entre os documentos eletrônicos abrangidos pelas orientações de 2026, NF-e, NFC-e, CT-e e outros documentos, cada qual sujeito aos leiautes técnicos aplicáveis. Ver \cite{RFB_ORIENT_2026}.
\end{fonte}

% ============================================================
\section{Especificação mínima recomendada para um MVP}
% ============================================================

A versão mais enxuta do projeto deve ser definida pela pergunta de cálculo. Se o cenário não exige memória de períodos anteriores, pode-se testar
\begin{equation}
\boxed{
\zeta_t^{\mathrm{MVP}}
=(u_t^{\min},\rho_t,\eta_t).
}
\label{eq:mvp}
\end{equation}
Se qualquer operador depender de saldos carregados, $x_t^{\min}$ volta a ser obrigatório.

\subsection{Dicionário mínimo de dados}

\renewcommand{\arraystretch}{1.15}
\small
\begin{longtable}{>{\raggedright\arraybackslash}p{2.0cm}>{\raggedright\arraybackslash}p{2.9cm}>{\raggedright\arraybackslash}p{4.1cm}>{\raggedright\arraybackslash}p{4.8cm}}
\toprule
\textbf{Objeto} & \textbf{Função no modelo} & \textbf{Fonte prática candidata} & \textbf{Critério de inclusão}\\
\midrule
\endfirsthead
\toprule
\textbf{Objeto} & \textbf{Função no modelo} & \textbf{Fonte prática candidata} & \textbf{Critério de inclusão}\\
\midrule
\endhead
$x_t^{\min}$ & estados com memória & ECD, ECF, controles fiscais/auxiliares & incluir se algum operador lê saldo de período anterior\\
$u_t^{\min}$ & operações necessárias ao cálculo & documentos fiscais e sistemas transacionais & projeção suficiente para $\cH$\\
$\rho_t$ & regime aplicável & cadastro fiscal, opção/regime declarado, documentação fiscal & obrigatório em cenários com regras dependentes de regime\\
$\eta_t$ & fatos da entidade & dados cadastrais e mestres & incluir atributos usados por admissibilidade ou regras\\
$\ell_{k,t}^{\min}$ & códigos fiscais/contábeis & documento fiscal, cadastro de item/serviço, escrituração & incluir códigos consumidos por $\cB_j,\cA_j,\cC_j$\\
$g_{k,t}$ & local da operação & documento fiscal/cadastro & incluir quando a regra dependa do local\\
$h_{k,t}$ & tratamento específico & classificação fiscal/regra aplicada & incluir quando houver regime específico/diferenciado\\
$\Prov(p)$ & auditoria normativa & repositório oficial de normas e documentação técnica & obrigatório para qualquer parâmetro codificado\\
\bottomrule
\end{longtable}
\normalsize

\begin{formalizacao}
O MVP não é definido por ``quantidade pequena de campos'', mas pela condição
\[
\text{informação coletada}
\;\Longrightarrow\;
\text{resultado identificável, reproduzível e normativamente rastreável}.
\]
\end{formalizacao}

% ============================================================
\section{Arquitetura consolidada do Volume II}
% ============================================================

A extensão completa pode ser resumida pelo fluxo
\begin{equation}
\boxed{
\cQ_t
\xrightarrow{\cN_t}
\bar\zeta_t
\xrightarrow{(\rho_t,\eta_t)}
\zeta_t
\xrightarrow{\mathfrak E_t(\Theta_t)}
\Theta_t^{\mathrm{eff}}
\xrightarrow{\cH}
Y_t
\xrightarrow{\Prov}
\text{resultado auditável}.
}
\end{equation}

O Volume I responde principalmente \emph{como operações viram demonstrações e tributos}. O Volume II acrescenta quatro perguntas operacionais:
\begin{enumerate}
  \item qual configuração tributária da entidade seleciona as regras relevantes;
  \item quais dados empresariais são suficientes para executar os operadores;
  \item qual grandeza de impacto será comparada no contrafactual;
  \item de qual dispositivo oficial veio cada parâmetro implementado.
\end{enumerate}

% ============================================================
\appendix
\section{Matriz de diferencial: conversa posterior versus Volume I}
% ============================================================

\begin{longtable}{p{3.6cm}p{4.1cm}p{6.1cm}}
\toprule
\textbf{Tema da conversa posterior} & \textbf{Situação no Volume I} & \textbf{Tratamento neste Volume II}\\
\midrule
\endfirsthead
\toprule
\textbf{Tema da conversa posterior} & \textbf{Situação no Volume I} & \textbf{Tratamento neste Volume II}\\
\midrule
\endhead
Regime tributário $\rho_t$ & ausente como objeto independente & definição por eixos, admissibilidade e seleção de regras efetivas\\
Características $\eta_t$ da entidade & implícitas em atributos e regras & vetor explícito de fatos de contexto\\
Vetor mínimo $z_t$ da conversa & não existente; $z_t$ já tinha outro significado & renomeado para $\zeta_t$ e definido por projeções suficientes\\
Plano de contas/Razão + documentos fiscais & fontes mencionadas em prosa & camada $\cQ_t$ e operador de normalização $\cN_t$\\
Simulador contrafactual & já existia para $\Theta^{old}/\Theta^{new}$ & decomposição entre efeito de lei, efeito de regime e interação\\
``Despesa com impostos'' & não separada como alvo de saída & distinção entre apuração, recolhimento, caixa, despesa contábil e crédito\\
Fontes oficiais & bibliografia e hierarquia já presentes & proveniência por parâmetro, versionamento e classes $\cS$\\
MVP de dados & ausente & critério de suficiência funcional e dicionário mínimo\\
\bottomrule
\end{longtable}

% ============================================================
\section{Símbolos adicionais}
% ============================================================

\begin{longtable}{p{3.1cm}p{3.0cm}p{7.4cm}}
\toprule
\textbf{Símbolo} & \textbf{Tipo} & \textbf{Definição}\\
\midrule
\endfirsthead
\toprule
\textbf{Símbolo} & \textbf{Tipo} & \textbf{Definição}\\
\midrule
\endhead
$\rho_t$ & configuração tributária & regime/enquadramento da entidade no período\\
$\eta_t$ & vetor factual & atributos da entidade relevantes à regra/admissibilidade\\
$\zeta_t$ & vetor de entrada & $(x_t^{\min},u_t^{\min},\rho_t,\eta_t)$\\
$\bar\zeta_t$ & base econômica & vetor de entrada sem a escolha de regime\\
$\mathfrak E_{j,t}$ & seletor & transforma regra geral em regra efetivamente aplicável\\
$\chi_t$ & indicador & verifica admissibilidade de um regime no cenário\\
$\mathsf R_t^{\mathrm{adm}}$ & conjunto & regimes admissíveis à entidade no cenário\\
$\cQ_t$ & pacote de fontes & ECD, ECF, documentos fiscais, cadastro e auxiliares\\
$\cN_t$ & operador ETL & normaliza fontes reais no esquema do simulador\\
$\Prov(p)$ & metadado & fonte, dispositivo, versão, vigência e consulta de $p$\\
$P_{j,t}^{\mathrm{cash}}$ & fluxo & pagamento financeiro do tributo no período\\
$E_{j,t}^{\mathrm{DRE}}$ & resultado & efeito reconhecido no resultado conforme norma aplicável\\
\bottomrule
\end{longtable}

% ============================================================
\section{Siglas adicionais}
% ============================================================

\begin{tabularx}{\textwidth}{>{\bfseries}p{2.6cm}X}
\toprule
ECD & Escrituração Contábil Digital\\
ECF & Escrituração Contábil Fiscal\\
SPED & Sistema Público de Escrituração Digital\\
RFB & Receita Federal do Brasil\\
CFC & Conselho Federal de Contabilidade\\
CGSN & Comitê Gestor do Simples Nacional\\
MEI & Microempreendedor Individual\\
NF-e & Nota Fiscal Eletrônica\\
NFC-e & Nota Fiscal de Consumidor Eletrônica\\
CT-e & Conhecimento de Transporte Eletrônico\\
MVP & Produto Mínimo Viável, no sentido de primeira implementação funcional\\
\bottomrule
\end{tabularx}

% ============================================================
\section{Referências normativas, contábeis e operacionais}
% ============================================================

\begin{thebibliography}{99}

\bibitem{LC214}
BRASIL. \textit{Lei Complementar nº 214, de 16 de janeiro de 2025, texto compilado}. Institui IBS, CBS e IS e disciplina, entre outros temas, o regime regular de IBS/CBS. Disponível em: \url{https://www.planalto.gov.br/ccivil_03/leis/lcp/lcp214compilado.htm}. Acesso em: 1 set. 2026.

\bibitem{RFB_RTC_LEG}
RECEITA FEDERAL DO BRASIL. \textit{Legislação da Reforma Tributária do Consumo}. Página atualizada em 18 ago. 2026, com EC 132/2023, LC 214/2025, LC 227/2026, Decreto 12.955/2026 e atos conjuntos/técnicos. Disponível em: \url{https://www.gov.br/receitafederal/pt-br/acesso-a-informacao/acoes-e-programas/programas-e-atividades/reforma-tributaria-do-consumo/legislacao}. Acesso em: 1 set. 2026.

\bibitem{RFB_SIMPLES_2026}
RECEITA FEDERAL DO BRASIL. \textit{CGSN atualiza regras do Simples Nacional para adequação à Reforma Tributária do Consumo}. 11 ago. 2026. Disponível em: \url{https://www.gov.br/receitafederal/pt-br/assuntos/noticias/2026/agosto/cgsn-atualiza-regras-do-simples-nacional-para-adequacao-a-reforma-tributaria-do-consumo/}. Acesso em: 1 set. 2026.

\bibitem{RFB_ECF_2026}
RECEITA FEDERAL DO BRASIL. \textit{ECF e Tributos da Pessoa Jurídica -- Perguntas e Respostas 2026}. Atualizado em 31 jul. 2026. Disponível em: \url{https://www.gov.br/receitafederal/pt-br/centrais-de-conteudo/publicacoes/perguntas-e-respostas/ecf}. Acesso em: 1 set. 2026.

\bibitem{ECD2026}
SISTEMA PÚBLICO DE ESCRITURAÇÃO DIGITAL. \textit{Manual de Orientação da ECD -- Leiaute 9 -- Atualização janeiro de 2026}. Anexo ao ADE Cofis nº 1/2026. Disponível em: \url{https://sped.rfb.gov.br/item/show/7990}. Acesso em: 1 set. 2026.

\bibitem{RFB_ORIENT_2026}
RECEITA FEDERAL DO BRASIL. \textit{Orientações da Reforma Tributária para 2026}. Atualizado em 6 maio 2026. Define, entre outras obrigações, emissão de documentos fiscais eletrônicos com CBS e IBS individualizados por operação conforme notas técnicas. Disponível em: \url{https://www.gov.br/receitafederal/pt-br/acesso-a-informacao/acoes-e-programas/programas-e-atividades/reforma-tributaria-do-consumo/orientacoes-2026}. Acesso em: 1 set. 2026.

\bibitem{RFB_CURSO}
RECEITA FEDERAL DO BRASIL; CONSELHO FEDERAL DE CONTABILIDADE. \textit{Curso Reforma Tributária do Consumo -- programação e materiais de apoio}. Em 2026, o Módulo 1 cobre fato gerador, base de cálculo, alíquotas, local da operação, não cumulatividade e apuração. Disponível em: \url{https://www.gov.br/receitafederal/pt-br/acesso-a-informacao/acoes-e-programas/programas-e-atividades/reforma-tributaria-do-consumo/curso/programacao}. Acesso em: 1 set. 2026.

\bibitem{CPC16}
COMITÊ DE PRONUNCIAMENTOS CONTÁBEIS. \textit{CPC 16 (R1) -- Estoques}. O custo de aquisição exclui tributos recuperáveis junto ao fisco. Sumário disponível em: \url{https://www.cpc.org.br/Arquivos/Documentos/244_CPC_16_Sumario.pdf}. Acesso em: 1 set. 2026.

\bibitem{CPC32}
COMITÊ DE PRONUNCIAMENTOS CONTÁBEIS. \textit{CPC 32 -- Tributos sobre o Lucro}. Define, entre outros objetos, despesa tributária corrente e diferida. Disponível em: \url{https://www.cpc.org.br/Arquivos/Documentos/340_CPC_32_rev%2013.pdf}. Acesso em: 1 set. 2026.

\end{thebibliography}

\end{document}
